% Options for packages loaded elsewhere
\PassOptionsToPackage{unicode}{hyperref}
\PassOptionsToPackage{hyphens}{url}
\documentclass[
  12pt,
]{article}
\usepackage{xcolor}
\usepackage[a4paper, margin=2.5cm]{geometry}
\usepackage{amsmath,amssymb}
\setcounter{secnumdepth}{-\maxdimen} % remove section numbering
\usepackage{iftex}
\ifPDFTeX
  \usepackage[T1]{fontenc}
  \usepackage[utf8]{inputenc}
  \usepackage{textcomp} % provide euro and other symbols
\else % if luatex or xetex
  \usepackage{unicode-math} % this also loads fontspec
  \defaultfontfeatures{Scale=MatchLowercase}
  \defaultfontfeatures[\rmfamily]{Ligatures=TeX,Scale=1}
\fi
\usepackage{lmodern}
\ifPDFTeX\else
  % xetex/luatex font selection
  \setmainfont[]{Georgia}
\fi
% Use upquote if available, for straight quotes in verbatim environments
\IfFileExists{upquote.sty}{\usepackage{upquote}}{}
\IfFileExists{microtype.sty}{% use microtype if available
  \usepackage[]{microtype}
  \UseMicrotypeSet[protrusion]{basicmath} % disable protrusion for tt fonts
}{}
\makeatletter
\@ifundefined{KOMAClassName}{% if non-KOMA class
  \IfFileExists{parskip.sty}{%
    \usepackage{parskip}
  }{% else
    \setlength{\parindent}{0pt}
    \setlength{\parskip}{6pt plus 2pt minus 1pt}}
}{% if KOMA class
  \KOMAoptions{parskip=half}}
\makeatother
\usepackage{longtable,booktabs,array}
\usepackage{calc} % for calculating minipage widths
% Correct order of tables after \paragraph or \subparagraph
\usepackage{etoolbox}
\makeatletter
\patchcmd\longtable{\par}{\if@noskipsec\mbox{}\fi\par}{}{}
\makeatother
% Allow footnotes in longtable head/foot
\IfFileExists{footnotehyper.sty}{\usepackage{footnotehyper}}{\usepackage{footnote}}
\makesavenoteenv{longtable}
\usepackage{graphicx}
\makeatletter
\newsavebox\pandoc@box
\newcommand*\pandocbounded[1]{% scales image to fit in text height/width
  \sbox\pandoc@box{#1}%
  \Gscale@div\@tempa{\textheight}{\dimexpr\ht\pandoc@box+\dp\pandoc@box\relax}%
  \Gscale@div\@tempb{\linewidth}{\wd\pandoc@box}%
  \ifdim\@tempb\p@<\@tempa\p@\let\@tempa\@tempb\fi% select the smaller of both
  \ifdim\@tempa\p@<\p@\scalebox{\@tempa}{\usebox\pandoc@box}%
  \else\usebox{\pandoc@box}%
  \fi%
}
% Set default figure placement to htbp
\def\fps@figure{htbp}
\makeatother
\setlength{\emergencystretch}{3em} % prevent overfull lines
\providecommand{\tightlist}{%
  \setlength{\itemsep}{0pt}\setlength{\parskip}{0pt}}
\usepackage{bookmark}
\IfFileExists{xurl.sty}{\usepackage{xurl}}{} % add URL line breaks if available
\urlstyle{same}

% Running margin label
\usepackage[useregional]{datetime2}
\usepackage{eso-pic}
\usepackage{tikz}
\usetikzlibrary{calc}
\newcommand{\mymarginlabel}{\textcopyright\ 2026 David Ewing dew59@uclive.ac.nz | \DTMnow\ | \jobname.tex}
\AddToShipoutPictureBG{%
  \begin{tikzpicture}[remember picture,overlay]
    \node[rotate=90, anchor=north]
      at ($(current page.north west)+(1.4cm,-13cm)$)
      {\scriptsize\ttfamily \mymarginlabel};
  \end{tikzpicture}%
}

\hypersetup{
  pdftitle={Tweet Sentiment Label Quality: A Text Classification Analysis},
  hidelinks,
  pdfcreator={LaTeX via pandoc}}

\title{Tweet Sentiment Label Quality: A Text Classification Analysis}
\author{}
\date{}

\begin{document}
\maketitle

\emph{David Ewing · Student ID: 82171165\\
DIGI405 Text Analysis · University of Canterbury · May 2026}

\section{1 Problem Statement}\label{problem-statement}

This project examines how well an automated text classifier can
reproduce the sentiment labels applied to a large sample of tweets. The
dataset is the Cardiff NLP tweet\_eval sentiment subset ---
45,615 tweets labelled as negative, neutral, or positive
by human annotators. Labelling was conducted as part of the SemEval 2017
Task 4A shared task on sentiment analysis in Twitter (Rosenthal et al.,
2017), and the labels are widely used as an NLP benchmark.

My central question is: how reliably can a text classifier reproduce
these labels, and what does classifier performance reveal about the
quality and consistency of the annotations? This matters because weak
classification performance may reflect not just the difficulty of the
task but genuine inconsistency in how the labels were assigned. A
classifier that performs poorly on neutral tweets may be revealing that
the boundary between neutral and mildly positive or mildly negative is
genuinely ambiguous --- that different annotators drew that line
differently. Put differently, low F1 may be evidence of noisy labels as
much as evidence of an inadequate model.

Three sub-questions guide the analysis:

\begin{enumerate}
\def\labelenumi{\arabic{enumi}.}
\item
  What kinds of problem are present in the labelling of tweets?
\item
  What types of tweet are systematically hard to classify?
\item
  Which features are most informative for distinguishing sentiment
  classes?
\item
  Does removing the neutral class substantially improve classifier
  performance, and what does this imply about the consistency of the
  neutral label?
\end{enumerate}

Macro average F1 is used as the primary evaluation metric. The
validation split (2,000 tweets) serves as the test set, following the
notebook template. Because the training data is class-imbalanced ---
neutral tweets account for roughly 45\% of the full training set --- the
random undersampling is applied to reduce all classes to the size of
the smallest (7,093 negative tweets), producing a balanced training set
of 21,279 examples. This approach was preferred over oversampling
because it avoids synthesising or duplicating examples that may
reinforce noisy labels. Macro F1 therefore averages equally across all
three classes, rather than being dominated by the majority class. This
is appropriate because the research question treats all three classes as
equally important: a metric that concentrates weight on neutral tweets
would obscure precisely the per-class performance differences this
analysis is designed to detect. Per-class precision and recall are also
reported. Accuracy is not used as the primary metric since it rewards
over-predicting the majority class.

\section{2 Dataset}\label{dataset}

The dataset is the sentiment subset of the Cardiff NLP tweet\_eval
benchmark (Barbieri et al., 2020). It was collected from Twitter and
annotated through the SemEval 2017 Task 4A competition, where teams
competed to classify tweet sentiment. Labels --- negative, neutral, and
positive --- were assigned by crowdsourced annotators, with final labels
reflecting majority vote across multiple annotations.

\emph{Table 1. Dataset split sizes and class distribution.}

\begin{longtable}[]{@{}
  >{\raggedright\arraybackslash}p{(\linewidth - 8\tabcolsep) * \real{0.2000}}
  >{\raggedright\arraybackslash}p{(\linewidth - 8\tabcolsep) * \real{0.2000}}
  >{\raggedright\arraybackslash}p{(\linewidth - 8\tabcolsep) * \real{0.2000}}
  >{\raggedright\arraybackslash}p{(\linewidth - 8\tabcolsep) * \real{0.2000}}
  >{\raggedright\arraybackslash}p{(\linewidth - 8\tabcolsep) * \real{0.2000}}@{}}
\toprule\noalign{}
\begin{minipage}[b]{\linewidth}\raggedright
\textbf{Split}
\end{minipage} & \begin{minipage}[b]{\linewidth}\raggedright
\textbf{Total}
\end{minipage} & \begin{minipage}[b]{\linewidth}\raggedright
\textbf{Negative}
\end{minipage} & \begin{minipage}[b]{\linewidth}\raggedright
\textbf{Neutral}
\end{minipage} & \begin{minipage}[b]{\linewidth}\raggedright
\textbf{Positive}
\end{minipage} \\
\midrule\noalign{}
\endhead
\bottomrule\noalign{}
\endlastfoot
Train (raw) & 45,615 & 7,093 & $\approx$19,800 & $\approx$18,700 \\
Train (after undersampling) & 21,279 & 7,093 & 7,093 & 7,093 \\
Validation / Test & 2,000 & 312 & 869 & 819 \\
Test (held out) & 12,284 & --- & --- & --- \\
\end{longtable}

After applying random undersampling (random\_state=82171165) to balance
the training classes, the training set is reduced to 21,279 tweets ---
7,093 per class. The validation split (2,000 tweets) is used as the
evaluation set. The test split (12,284 tweets) is held out and not used
in this report. Exact per-class counts at each stage of processing are
shown in notebook §2.5.

The class distribution in the validation set is unbalanced: 312
negative, 869 neutral, and 819 positive tweets. This asymmetry --- far
fewer negative examples than neutral or positive --- means per-class
results need careful interpretation. Poor performance on negative tweets
may partly reflect fewer test examples as well as weaker training
signal.

Tweets are short, informal, and noisy: hashtags, @-mentions, URLs,
slang, emoji, and abbreviations are common. They are also highly
context-dependent --- the same words can carry different sentiment
depending on the event being discussed or who is being addressed. These
properties make sentiment classification in Twitter data intrinsically
harder than in product reviews or news text.

\section{3 Model}\label{model}

Two classifiers are evaluated: logistic regression and a shallow
decision tree. Both use the same scikit-learn pipeline (notebook section
3), fitted to the balanced training set and evaluated on the validation
split.

The pipeline has four steps. First, a TextCleaner component strips
whitespace. Second, a SpacyPreprocessor tokenises each tweet using the
en\_core\_web\_sm spaCy model and caches the results in an SQLite
feature store (assignment-sentiment.sqlite). Third, a FeatureUnion
assembles the feature matrix. In this initial run the only active
feature block is a TokensVectorizer: unigram token counts for the top
100 most frequent tokens by document frequency, after removing stop
words and punctuation, with lowercasing applied. Fourth, the feature
matrix is passed to the classifier.

Logistic regression is configured with max\_iter=5000, random\_state=42.
The decision tree uses max\_depth=3, random\_state=42. The shallow depth
restricts the tree to at most seven leaf nodes, forcing selection of
only the most discriminating tokens and producing a human-interpretable
model. Both classifiers share identical preprocessing, achieved by
calling set\_params(classifier=clf) inside a loop over both
configurations.

\textbf{Dummy-image}

\emph{Figure 1. Scikit-learn pipeline.}

This run uses a minimal, interpretable baseline feature set. The model
plan lists additional feature types to explore in subsequent runs:
TF-IDF weighted unigrams with a larger vocabulary, bigrams, VADER
sentiment lexicon counts, document-level text statistics (tweet length,
punctuation counts), and POS tag sequences. A binary task (negative vs
positive only) is also planned.

\section{4 Results}\label{results}

Tables 2 and 3 show the classification report for logistic regression
and decision tree respectively, on the 2,000-tweet validation split.
Figures 2 and 3 show the corresponding confusion matrices. Figure 4
shows the top LR feature coefficients per class. Figure 5 shows the
decision tree structure.

\emph{Table 2. Logistic Regression --- classification report (unigram
counts, top 100 features).}

\begin{longtable}[]{@{}
  >{\raggedright\arraybackslash}p{(\linewidth - 8\tabcolsep) * \real{0.2000}}
  >{\raggedright\arraybackslash}p{(\linewidth - 8\tabcolsep) * \real{0.2000}}
  >{\raggedright\arraybackslash}p{(\linewidth - 8\tabcolsep) * \real{0.2000}}
  >{\raggedright\arraybackslash}p{(\linewidth - 8\tabcolsep) * \real{0.2000}}
  >{\raggedright\arraybackslash}p{(\linewidth - 8\tabcolsep) * \real{0.2000}}@{}}
\toprule\noalign{}
\begin{minipage}[b]{\linewidth}\raggedright
\textbf{Class}
\end{minipage} & \begin{minipage}[b]{\linewidth}\raggedright
\textbf{Precision}
\end{minipage} & \begin{minipage}[b]{\linewidth}\raggedright
\textbf{Recall}
\end{minipage} & \begin{minipage}[b]{\linewidth}\raggedright
\textbf{F1}
\end{minipage} & \begin{minipage}[b]{\linewidth}\raggedright
\textbf{Support}
\end{minipage} \\
\midrule\noalign{}
\endhead
\bottomrule\noalign{}
\endlastfoot
negative & 0.245 & 0.446 & 0.316 & 312 \\
neutral & 0.517 & 0.534 & 0.525 & 869 \\
positive & 0.631 & 0.413 & 0.499 & 819 \\
macro avg & 0.464 & 0.464 & 0.447 & 2000 \\
accuracy & --- & --- & 0.470 & 2000 \\
\end{longtable}

\emph{Table 3. Decision Tree (max\_depth=3) --- classification report
(unigram counts, top 100 features).}

\begin{longtable}[]{@{}
  >{\raggedright\arraybackslash}p{(\linewidth - 8\tabcolsep) * \real{0.2000}}
  >{\raggedright\arraybackslash}p{(\linewidth - 8\tabcolsep) * \real{0.2000}}
  >{\raggedright\arraybackslash}p{(\linewidth - 8\tabcolsep) * \real{0.2000}}
  >{\raggedright\arraybackslash}p{(\linewidth - 8\tabcolsep) * \real{0.2000}}
  >{\raggedright\arraybackslash}p{(\linewidth - 8\tabcolsep) * \real{0.2000}}@{}}
\toprule\noalign{}
\begin{minipage}[b]{\linewidth}\raggedright
\textbf{Class}
\end{minipage} & \begin{minipage}[b]{\linewidth}\raggedright
\textbf{Precision}
\end{minipage} & \begin{minipage}[b]{\linewidth}\raggedright
\textbf{Recall}
\end{minipage} & \begin{minipage}[b]{\linewidth}\raggedright
\textbf{F1}
\end{minipage} & \begin{minipage}[b]{\linewidth}\raggedright
\textbf{Support}
\end{minipage} \\
\midrule\noalign{}
\endhead
\bottomrule\noalign{}
\endlastfoot
negative & 0.000 & 0.000 & 0.000 & 312 \\
neutral & 0.454 & 0.978 & 0.620 & 869 \\
positive & 0.798 & 0.121 & 0.210 & 819 \\
macro avg & 0.417 & 0.366 & 0.277 & 2000 \\
accuracy & --- & --- & 0.474 & 2000 \\
\end{longtable}

\emph{Table 4. Macro F1 comparison: Logistic Regression vs Decision
Tree.}

\begin{longtable}[]{@{}
  >{\raggedright\arraybackslash}p{(\linewidth - 6\tabcolsep) * \real{0.2500}}
  >{\raggedright\arraybackslash}p{(\linewidth - 6\tabcolsep) * \real{0.2500}}
  >{\raggedright\arraybackslash}p{(\linewidth - 6\tabcolsep) * \real{0.2500}}
  >{\raggedright\arraybackslash}p{(\linewidth - 6\tabcolsep) * \real{0.2500}}@{}}
\toprule\noalign{}
\begin{minipage}[b]{\linewidth}\raggedright
\textbf{Class}
\end{minipage} & \begin{minipage}[b]{\linewidth}\raggedright
\textbf{LR F1}
\end{minipage} & \begin{minipage}[b]{\linewidth}\raggedright
\textbf{DT F1}
\end{minipage} & \begin{minipage}[b]{\linewidth}\raggedright
\textbf{LR − DT}
\end{minipage} \\
\midrule\noalign{}
\endhead
\bottomrule\noalign{}
\endlastfoot
negative & 0.316 & 0.000 & +0.316 \\
neutral & 0.525 & 0.620 & −0.095 \\
positive & 0.499 & 0.210 & +0.289 \\
macro & 0.447 & 0.277 & +0.170 \\
\end{longtable}

\emph{Figure 2. Logistic Regression --- confusion matrix.}

\emph{Figure 3. Decision Tree --- confusion matrix.}

\emph{Figure 4. Logistic Regression --- top 20 feature coefficients per
class.}

\emph{Figure 5. Decision Tree visualisation (max\_depth=3).}

\section{5 Discussion}\label{discussion}

\subsection{5.1 Overall performance}\label{overall-performance}

The logistic regression baseline achieved a macro F1 of 0.447. This is
below 0.5, a rough reference point for a classifier that handles all
three classes with at least some reliability. The result is consistent
with the known difficulty of three-class sentiment classification in
tweets: the task is harder than binary sentiment in longer, more formal
text, and the class boundaries in Twitter data are inherently fuzzy.

The decision tree achieved a lower macro F1 of 0.277, despite marginally
higher accuracy (0.474 vs 0.470). The accuracy paradox is instructive:
the test set contains 869 neutral tweets (43.5\%), so a classifier that
predicts neutral for almost all inputs achieves near-50\% accuracy while
producing a very low macro F1. This confirms that macro F1 is the
correct metric for this dataset.

\subsection{5.2 Class-level analysis}\label{class-level-analysis}

The most striking pattern in the LR results is the asymmetry across
classes. Positive tweets achieved the highest F1 (0.499), driven by high
precision (0.631) but only moderate recall (0.413): the model is
conservative about predicting positive, but when it does it is often
right. Negative tweets achieved the lowest F1 (0.316), with very low
precision (0.245) but moderate recall (0.446): the model over-predicts
negative, and most tweets predicted as negative were actually neutral or
positive. Neutral sat in the middle (F1=0.525).

The fact that neutral is not the worst-performing class is somewhat
surprising. The conventional view --- supported by the assignment notes
and by Kenyon-Dean et al. (2018) --- is that neutral is the hardest
class because it lacks distinctive lexical markers. The LR results do
not clearly support this: neutral F1 (0.525) exceeds negative F1
(0.316). This may be because the neutral class is the largest in the
test set (869 examples), giving the model more evaluation signal, or
because neutral tweets tend to contain common, non-sentiment-loaded
words that are reliably represented in the top-100 unigram vocabulary.
Negative tweets may use more varied and context-dependent language, and
the test set has only 312 negative examples --- the fewest of the three
classes.

The decision tree result tells a clearer story. With max\_depth=3 and
100 unigram count features, the tree predicted neutral for almost all
inputs (recall=0.978 for neutral, near-zero for both other classes).
This is consistent with a shallow tree learning a dominant
majority-class heuristic --- even though training data was balanced, the
learned split rules apparently generalise better for neutral than for
the other classes.

\subsection{5.3 What does this tell us about label
quality?}\label{what-does-this-tell-us-about-label-quality}

The central question of this project is not simply "how good is the
classifier?" but "what does classifier performance tell us about the
labels?" The low macro F1 is partly attributable to model limitations
--- the feature set is minimal --- but it is also consistent with
genuine label noise in the dataset.

The most suggestive evidence comes from the negative class. Low
precision (0.245) means the model is making many false positive
predictions for negative sentiment. From a labelling perspective: if the
model predicts negative based on reasonable lexical cues but the gold
label is neutral, this may indicate cases where annotators labelled a
borderline tweet as neutral rather than negative. The SemEval annotation
guidelines used neutral as a default for tweets that did not clearly
express positive or negative sentiment --- which may have produced a
diffuse, underspecified neutral category that overlaps substantially
with mild negative expression.

Kenyon-Dean et al. (2018) argue that sentiment classification in tweets
is "complicated" precisely because sentiment is a property of the
author\textquotesingle s relationship to the content, not just of the
words used. The same phrase ("not bad at all") may be positive or ironic
depending on context. Unigram count features, which discard word order,
negation scope, and context, are poorly equipped to capture these
distinctions. The weak results are therefore expected --- but they also
suggest that the hard cases may involve genuine ambiguity that even
human annotators would disagree on.

\subsection{5.4 Limitations}\label{limitations}

This analysis is preliminary. Only one feature configuration has been
evaluated --- 100 unigram count features --- and no comparison across
feature types has been made. The planned additional experiments (VADER
features, binary task, bigrams, text statistics) are not yet complete.
The decision tree result cannot be fully interpreted without examining
the tree visualisation (Figure 5, currently a placeholder). The
misclassification analysis in notebook section 4.2 was run but
individual tweet examples have not yet been examined in detail.

The binary task experiment (negative vs positive only) is needed to
determine whether neutral is primarily responsible for the overall weak
performance. Until this is run, the relative contribution of the neutral
class cannot be quantified.

The analysis is also constrained to the validation split. The held-out
test split (12,284 tweets) has not been used and would provide a more
reliable estimate of generalisation performance.

\section{References}\label{references}

Barbieri, F., Camacho-Collados, J., Espinosa-Anke, L., \& Neves, L.
(2020). TweetEval: Unified Benchmark and Comparative Evaluation for
Tweet Classification. In Findings of the Association for Computational
Linguistics: EMNLP 2020 (pp. 1644--1650).
https://doi.org/10.18653/v1/2020.findings-emnlp.148

Hutto, C. J., \& Gilbert, E. (2014). VADER: A Parsimonious Rule-Based
Model for Sentiment Analysis of Social Media Text. In Proceedings of the
Eighth International AAAI Conference on Weblogs and Social Media (pp.
216--225).

Kenyon-Dean, K., Ahmed, E., Bhosale, S., Brooks, B., Laframboise, C.,
Roper, F., ... \& Singh, S. (2018). Sentiment Analysis:
It\textquotesingle s Complicated! In Proceedings of the 2018 Conference
of the North American Chapter of the Association for Computational
Linguistics: Human Language Technologies, Volume 1 (pp. 1886--1895).
https://doi.org/10.18653/v1/N18-1171

Rosenthal, S., Farra, N., \& Nakov, P. (2017). SemEval-2017 Task 4:
Sentiment Analysis in Twitter. In Proceedings of the 11th International
Workshop on Semantic Evaluation (SemEval-2017) (pp. 502--518).
\url{https://doi.org/10.18653/v1/S17-2088}

\section[Appendix\\
\strut \\
]{\texorpdfstring{Appendix\\
\strut \\
\protect\includegraphics[width=6.26806in,height=5.86181in]{./media/image1.png}}{Appendix  }}\label{appendix}

\includegraphics[width=6.26806in,height=7.86597in]{./media/image2.png}

\includegraphics[width=6.26806in,height=5.82431in]{./media/image3.png}

\end{document}
